% ===========================================================================
%  HaloNet — IEEE TIP Draft
%  Introduction + Related Work
% ===========================================================================

\section{Introduction}\label{sec:intro}

Infrared small target detection (IRSTD) plays an indispensable role in maritime surveillance, early-warning systems, and remote sensing~\cite{dnanet,uiunet,alcnet}.
Unlike generic object detection, infrared small targets typically occupy fewer than $9\times9$ pixels, lack discernible texture or shape cues, and are deeply embedded in spatially correlated background clutter~\cite{irstd_survey}.
Despite the remarkable progress brought by deep learning~\cite{dnanet,uiunet,sctransnet}, fully supervised methods remain fundamentally bottlenecked by the cost of acquiring pixel-level mask annotations.
Precisely delineating a $3\times3$ target against a noisy infrared background demands domain expertise and meticulous per-pixel labeling---a process that scales poorly to the large-scale datasets required for robust generalization.

To mitigate this annotation burden, recent works have explored \emph{weakly supervised} paradigms for IRSTD.
Notably, LESPS~\cite{lesps} introduces single-point supervision with a label evolution mechanism that iteratively refines pseudo-masks during training, while PAL~\cite{pal} further incorporates an active learning loop to progressively select the most informative annotation points.
Although these methods significantly reduce annotation effort, they share two fundamental limitations that hinder their deployment in operational scenarios.
First, \emph{single-point annotation in dense clutter is inherently fragile}: accurately clicking the geometric centroid of a dim, sub-$5$-pixel target embedded in heavy sea-surface or cloud clutter is both time-consuming and prone to human jittering error, as illustrated in Fig.~\ref{fig:annotation_comparison}.
Second, the iterative label evolution and active learning loops employed by LESPS and PAL introduce substantial computational overhead and training complexity, undermining the simplicity desirable for real-world deployment.

In contrast, \emph{bounding box} annotations represent the most natural and readily available form of weak supervision in practice.
Upstream early-warning systems---ranging from radar cueing to lightweight YOLO-based detectors~\cite{yolo}---routinely produce coarse bounding boxes as their primary output.
Requiring these operational pipelines to additionally supply pixel-level masks or geometrically precise center points is neither realistic nor cost-effective.
However, leveraging box annotations for IRSTD poses a unique and severe challenge: for a typical small target occupying fewer than $0.5\%$ of the box area, over $99\%$ of the enclosed pixels belong to background clutter.
Naively treating all box-interior pixels as foreground leads to catastrophic representation collapse, where the network learns to predict the entire box region rather than the compact target.

% --- Core Contribution Paragraph ---

In this paper, we propose \textbf{HaloNet}\footnote{The name alludes to the faint ``halo'' of physical signal surrounding an infrared point source.}, a \emph{physics-anchored pseudo-label generation} framework that converts coarse bounding boxes into high-fidelity pixel-level soft masks, \emph{without any learnable parameters, iterative refinement, or network architecture modification}.
Our key insight is that the radiometric properties of infrared imagery provide a principled statistical basis for distinguishing target pixels from background within a bounding box.
Concretely, we make two contributions at the pseudo-label generation level:

\begin{enumerate}
    \item \textbf{Background-SNR Posterior Mask (B-SNR).}
    Under a local Gaussian background assumption, we show that the per-pixel posterior probability of belonging to the target class admits a closed-form sigmoid expression:
    $W(i) = \sigma\bigl(\tau \cdot ({I(i) - \mu_\mathrm{ctx}})/{\sigma_\mathrm{ctx}}\bigr)$,
    where $\mu_\mathrm{ctx}$ and $\sigma_\mathrm{ctx}$ are the local background statistics estimated from the box context, and $\tau$ controls the decision sharpness.
    This formulation bridges classical signal-to-noise ratio (SNR) analysis with Bayesian posterior estimation, providing a statistically grounded binary soft mask.

    \item \textbf{Physics-Anchored Gaussian Refinement (PAG).}
    To further suppress residual clutter within the B-SNR mask, we exploit the physical point-spread characteristic of infrared targets.
    Using the full-width-at-half-maximum (FWHM) of the B-SNR response as an anchor, we perform an SNR-weighted spatial covariance estimation to derive an anisotropic Gaussian spatial prior.
    The resulting PAG mask smoothly concentrates supervision on the target's energy peak while gracefully attenuating toward the periphery, yielding a soft label that better reflects the underlying point-spread function.
\end{enumerate}

Beyond the pseudo-label generation itself, we provide a \textbf{theoretical analysis of the variance corruption phenomenon} that occurs when targets are extremely small relative to the bounding box.
Specifically, we derive a closed-form expression showing that the estimated background variance is inflated by a term proportional to $\alpha(1-\alpha)\Delta\mu^2$, where $\alpha$ denotes the target-to-box area ratio and $\Delta\mu$ the target-background intensity contrast.
This result rigorously explains why physics-based pseudo-labels degrade under extreme scale conditions (e.g., sub-$3\times3$ targets in NUDT-SIRST), establishing a principled failure boundary rather than concealing the limitation.

We summarize the main contributions of this work as follows:

\begin{itemize}
    \item We present \textbf{HaloNet}, the first box-supervised framework for infrared small target detection.
    By leveraging the radiometric statistics of infrared imagery, HaloNet generates high-quality pixel-level pseudo-masks from coarse bounding boxes in a single forward pass, with zero learnable parameters and no iterative refinement.

    \item We provide a \textbf{rigorous statistical foundation} for the proposed pseudo-label generation pipeline.
    The B-SNR mask is derived as a Bayesian posterior under a Gaussian background shift model, and the PAG refinement is formulated as an SNR-weighted spatial covariance estimation.
    Furthermore, we derive a \textbf{variance corruption bound} that quantitatively characterizes the theoretical degradation regime of physics-based methods under extreme target-to-box ratios.

    \item We conduct \textbf{extensive model-agnostic experiments} across three standard IRSTD benchmarks (NUAA-SIRST, NUDT-SIRST, IRSTD-1K) and three architecturally diverse backbones (DNANet, ACM, ALCNet).
    HaloNet consistently outperforms prior box-based baselines, achieves competitive or superior performance relative to single-point methods, and demonstrates strong robustness to bounding box perturbations.
\end{itemize}


% ===========================================================================
\section{Related Work}\label{sec:related}
% ===========================================================================

\subsection{Deep Learning for Infrared Small Target Detection}\label{sec:related_dl}

The past five years have witnessed rapid architectural innovation in fully supervised IRSTD.
DNANet~\cite{dnanet} introduced dense nested connections within a U-Net++ encoder--decoder topology, coupled with a channel--spatial attention mechanism (Res-CBAM), establishing state-of-the-art segmentation performance on NUDT-SIRST.
UIU-Net~\cite{uiunet} proposed a nested ``U-Net in U-Net'' structure to capture both macro-level contextual semantics and micro-level local details for small targets.
ACM~\cite{acm} designed an asymmetric contextual modulation module to fuse multi-scale features for improved target-background discrimination, while ALCNet~\cite{alcnet} further incorporated local contrast attention into a Feature Pyramid Network decoder.
More recently, SCTransNet~\cite{sctransnet} and ILNet~\cite{ilnet} have explored Transformer-based and low-level feature enhancement strategies, respectively.

Despite their strong performance, these methods share a common and fundamental limitation: they are critically dependent on high-quality pixel-level mask annotations.
As the community scales toward larger and more diverse datasets, the annotation bottleneck becomes the dominant obstacle, motivating the exploration of weaker forms of supervision.

\subsection{Weakly Supervised Infrared Small Target Detection}\label{sec:related_ws}

Weakly supervised IRSTD has emerged as a promising direction to reduce annotation costs.
LESPS~\cite{lesps} pioneered single-point supervision for IRSTD, proposing a mapping degeneration framework in which pseudo-masks are iteratively evolved during network training via a label update mechanism driven by the network's own predictions.
PAL~\cite{pal} extended this paradigm by introducing a progressive active learning framework that strategically selects the most informative annotation points, achieving strong performance with minimal human intervention.

While these methods represent significant advances, they exhibit two structural limitations.
\emph{First}, both methods rely on multi-round iterative procedures---LESPS performs label evolution at every training epoch, and PAL requires multiple active learning cycles with human-in-the-loop annotation.
These iterative mechanisms increase training complexity and introduce sensitivity to the convergence behavior of the self-training loop.
\emph{Second}, the single-point annotation paradigm itself is fragile in operational settings: for extremely dim targets (SCR $<$ 2) in heavy clutter, reliably identifying and clicking the target centroid demands significant expertise and concentration, and the resulting labels are susceptible to spatial jittering that can mislead label evolution.

In contrast, our work departs from the single-point paradigm entirely, proposing instead to exploit bounding box annotations---a supervision signal that is naturally produced by upstream detection systems and requires no sub-pixel human precision.

\subsection{Box-Supervised Semantic Segmentation}\label{sec:related_box}

Box-supervised segmentation has been extensively studied in the natural image domain.
Early approaches such as BoxSup~\cite{boxsup} and SDI~\cite{sdi} used bounding boxes to generate coarse pseudo-masks via GrabCut~\cite{grabcut} or MCG~\cite{mcg} proposals, followed by iterative retraining.
More recent methods leverage class activation maps (CAMs)~\cite{cam} or affinity-based propagation~\cite{affinitynet} to refine box-level supervision into pixel-level pseudo-labels.
BoxInst~\cite{boxinst} and Box2Mask~\cite{box2mask} further advanced this line by incorporating projection-based losses and pairwise affinity constraints.

However, these methods are designed for natural images where targets exhibit rich texture, distinctive color, and clear semantic boundaries---cues that are entirely absent in infrared small target imagery.
An infrared point target is a near-isotropic energy blob spanning a few pixels, surrounded by spatially correlated clutter with similar intensity distributions.
Standard CAM-based or affinity-based approaches, which rely on semantic feature contrast to distinguish foreground from background, fail catastrophically in this regime because the network features of a sub-$5$-pixel target are indistinguishable from background noise at typical feature map resolutions.

This fundamental mismatch motivates our \emph{physics-driven} approach: rather than relying on learned semantic features to parse the box interior, we exploit the known radiometric properties of infrared point sources---specifically, the local signal-to-noise ratio and the physical point-spread function---to derive pseudo-masks that are grounded in the imaging physics rather than network-dependent heuristics.
